\SourceLeaf{151}%
\[
\begin{array}{c|lll}
\text{pro complexu}&\multicolumn{3}{c}{\text{positivi esse debent numeri}}\\\hline
C&\eta,&\xi-\eta,&p-\xi-\eta\\
C'&p-\xi,&\xi-\eta,&\xi+\eta-p\\
C''&p-\eta,&\eta-\xi,&\xi+\eta-p\\
C'''&\xi,&\eta-\xi,&p-\xi-\eta
\end{array}
\]

Ceterum vel nobis non monentibus quisque facile intelliget, in repraesentatione figurata numerorum complexorum (vid. art.~39) numeros systematis $S$ intra quadratum contineri, cuius latera iungant puncta numeros $0,a+bi,(1+i)(a+bi),i(a+bi)$ repraesentantia, et subdivisionem systematis $S$ respondere partitioni quadrati per rectas diagonales. Sed hocce loco ratiocinationibus pure arithmeticis uti maluimus, illustrationem per intuitionem figuratam lectori perito brevitatis caussa linquentes.

\Article{70}
Si quatuor numeri complexi $r=x+yi,r'=x'+y'i,r''=x''+y''i,r'''=x'''+y'''i$ ita inter se nexi sunt, ut habeatur $r'=m+ir,r''=m+ir'=(1+i)m-r,r'''=m+ir''=im-ir$, atque primus $r$ ad complexum $C$ pertinere supponitur, reliqui $r',r'',r'''$ resp. ad complexus $C',C'',C'''$ pertinebunt. Statuendo enim $\xi=ax+by,\eta=ay-bx,\xi'=ax'+by',\eta'=ay'-bx',\xi''=ax''+by'',\eta''=ay''-bx'',\xi'''=ax'''+by''',\eta'''=ay'''-bx'''$, invenitur
\[
\begin{aligned}
\eta&=p-\xi'=p-\eta''=\xi'''\\
\xi-\eta&=\xi'+\eta'-p=\eta''-\xi''=p-\xi'''-\eta'''\\
p-\xi-\eta&=\xi'-\eta'=\xi''+\eta''-p=\eta'''-\xi'''
\end{aligned}
\]
unde adiumento criteriorum theorematis veritas sponte demanat. Et quum rursus fiat $r=m+ir'''$, facile perspicitur, si $r$ supponatur pertinere ad $C'$, numeros $r',r'',r'''$ pertinere resp. ad $C'',C''',C$; si ille ad $C''$, hos ad $C''',C,C'$; denique si ille ad $C'''$, hos ad $C,C',C''$.

Simul hinc colligitur, in singulis complexibus $C,C',C'',C'''$ aeque multos numeros reperiri, puta $\frac14(p-1)$.

\par\medskip\Needspace{4\baselineskip}\pdfbookmark[1]{Art. 37}{paper-five-native151-art37}\centerline{\textbf{37.}}\nobreak\smallskip
\textsc{Theorema.} \emph{Si denotante $k$ integrum per $m$ non divisibilem singuli numeri complexus $C$ per $k$ multiplicantur, productorumque residuis simpliciter minimis secundum
\SourceLeaf{152}%
modulum $m$ inter complexus $C,C',C'',C'''$ distributis, multitudo eorum, quae ad singulos hos complexus pertinent, resp. per $c,c',c'',c'''$ denotatur: character numeri $k$ respectu moduli $m$ erit $\equiv c'+2c''+3c'''\pmod4$.}

\emph{Demonstr.} Sint illa $c$ residua minima ad $C$ pertinentia $\alpha,\beta,\gamma,\delta$ etc.; dein $c'$ residua ad $C'$ pertinentia haec $m+i\alpha',m+i\beta',m+i\gamma',m+i\delta'$ etc.; porro $c''$ residua ad $C''$ pertinentia haec $(1+i)m-\alpha'',(1+i)m-\beta'',(1+i)m-\gamma'',(1+i)m-\delta''$ etc.; denique $c'''$ residua ad $C'''$ pertinentia haec $im-i\alpha''',im-i\beta''',im-i\gamma''',im-i\delta'''$ etc. Iam consideremus quatuor producta, scilicet
\begin{enumerate}
\renewcommand{\labelenumi}{\arabic{enumi})}
\item productum ex omnibus $\frac14(p-1)$ numeris complexum $C$ constituentibus:
\item productum productorum, quae e multiplicatione singulorum horum numerorum per $k$ orta erant;
\item productum e residuis minimis horum productorum, puta e numeris $\alpha,\beta,\gamma,\delta$ etc., $m+i\alpha',m+i\beta'$ etc. etc.
\item productum ex omnibus $c+c'+c''+c'''$ numeris $\alpha,\beta,\gamma,\delta$ etc., $\alpha',\beta',\gamma',\delta'$ etc., $\alpha'',\beta'',\gamma'',\delta''$ etc.
\end{enumerate}
Denotando haec quatuor producta ordine suo per $P,P',P'',P'''$, manifesto erit
\[
\begin{gathered}
P'=k^{\frac14(p-1)}P,\quad P'\equiv P'',\quad
P''\equiv P'''i^{c'+2c''+3c'''}\pmod m
\end{gathered}
\]
et proin
\[
Pk^{\frac14(p-1)}\equiv P'''i^{c'+2c''+3c'''}\pmod m
\]

At facile perspicitur, numeros $\alpha',\allowbreak\beta',\allowbreak\gamma',\allowbreak\delta'$ etc., $\alpha'',\allowbreak\beta'',\allowbreak\gamma'',\allowbreak\delta''$ etc., $\alpha''',\allowbreak\beta''',\allowbreak\gamma''',\allowbreak\delta'''$ etc. omnes ad complexum $C$ pertinere, atque tum inter se tum a numeris $\alpha,\allowbreak\beta,\allowbreak\gamma,\allowbreak\delta$ etc. diversos esse, sicuti hi ipsi inter se diversi sint. Omnes itaque hi numeri simul sumti, et abstrahendo ab ordine, prorsus identici esse debent cum omnibus numeris complexum $C$ constituentibus, unde colligimus $P=P'''$, adeoque
\[
Pk^{\frac14(p-1)}\equiv Pi^{c'+2c''+3c'''}\pmod m
\]
Denique quum singuli factores producti $P$ per $m$ non sint divisibiles, hinc concluditur
\[
k^{\frac14(p-1)}\equiv i^{c'+2c''+3c'''}\pmod m
\]
unde $c'+2c''+3c'''$ erit character numeri $k$ respectu moduli $m$. Q.~E.~D.

\SourceLeaf{153}%
\Article{72}
Quo theorema generale art. praec. ad numerum $1+i$ applicari possit, complexum $C$ denuo in duos complexus minores $G$ et $G'$ subdividere oportet, et quidem referemus in complexum $G$ numeros eos $x+yi$, pro quibus $ax+by=\xi$ minor est quam $\frac12p$, in alterum $G'$ eos, pro quibus $\xi$ est maior quam $\frac12p$; multitudinem numerorum in complexibus $G,G'$ contentorum resp. per $g,g'$ denotabimus, unde erit $g+g'=\frac14(p-1)$.

Criterium completum numerorum ad $G$ pertinentium itaque erit, ut tres numeri $\eta,\xi-\eta,p-2\xi$ sint positivi: nam conditio tertia pro complexu $C$, secundum quam $p-\xi-\eta$ positivus esse debet, sub illis implicite iam continetur, quum sit $p-\xi-\eta=(\xi-\eta)+(p-2\xi)$. Perinde criterium completum numerorum ad $G'$ pertinentium consistet in valoribus positivis trium numerorum $\eta,p-\xi-\eta,2\xi-p$.

Hinc facile concluditur, productum cuiusvis numeri complexus $G$ per numerum $1+i$ pertinere ad complexum $C'''$; si enim statuitur
\[
\begin{gathered}
(x+yi)(1+i)=x'+y'i,\quad\text{atque}\quad ax'+by'=\xi',\ ay'-bx'=\eta',\ \text{invenitur}\\
\xi'=\xi-\eta,\quad\eta'-\xi'=2\eta,\quad p-\xi'-\eta'=p-2\xi
\end{gathered}
\]
i.~e. criterium pro numero $x+yi$ complexui $G$ subdito identicum est cum criterio pro numero $x'+y'i$ ad complexum $C'''$ pertinente.

Prorsus simili modo ostenditur, productum cuiusvis numeri complexus $G'$ per $1+i$ pertinere ad complexum $C''$.

Erit itaque, si in art. praec. ipsi $k$ valorem $1+i$ tribuimus, $c=0,c'=0,c''=g',c'''=g$, et proin character numeri $1+i$ fiet $3g+2g'=\frac12(p-1)+g$. Et quum characteres numerorum $i,-1$, sint $\frac14(p-1),\frac12(p-1)$, characteres numerorum $-1+i,-1-i,1-i$ resp. erunt $\frac34(p-1)+g,g,\frac14(p-1)+g$. Totus igitur rei cardo iam in investigatione numeri $g$ vertitur.

\Article{73}
Quae in artt.~69---72 exposuimus, proprie independentia sunt a suppositione, $m$ esse numerum primarium: abhinc vero saltem supponemus, $a$ imparem, $b$ parem esse, praetereaque $a,b$ et $a-b$ esse numeros positivos. Ante omnia limites valorum ipsius $x$ in complexu $G$ stabilire oportet.

\SourceLeaf{154}%
Statuendo $ay-bx=\eta,(a+b)x-(a-b)y=\zeta,p-2ax-2by=\theta$, criterium numerorum $x+yi$ ad complexum $G$ pertinentium consistit in tribus conditionibus, ut $\eta,\zeta,\theta$ sint numeri positivi. Quum fiat $px=(a-b)\eta+a\zeta,p(a-2x)=a\theta+2b\eta$, manifestum est, $x$ et $2a-x$ esse debere numeros positivos, sive $x$ alicui numerorum $1,2,3\ldots\frac12(a-1)$ aequalem. Porro quum sit $(a-b)\theta=2b\zeta+p(a-b-2x)$, patet, quamdiu $x$ minor sit quam $\frac12(a-b)$, conditionem secundam (iuxta quam $\zeta$ positivus esse debet) iam implicare tertiam (quod $\theta$ debet esse positivus); contra quoties $x$ sit maior quam $\frac12(a-b)$, conditionem secundam iam contineri sub tertia. Quamobrem pro valoribus ipsius $x$ his $1,2,3\ldots\frac12(a-b-1)$ tantummodo prospiciendum est, ut $\eta$ et $\zeta$ positivi evadant, sive ut $y$ maior sit quam $\frac{bx}{a}$ et minor quam $\frac{(a+b)x}{a-b}$: pro valore itaque tali dato ipsius $x$ aderunt numeri $x+yi$ omnino
\[
\left[\frac{(a+b)x}{a-b}\right]-\left[\frac{bx}{a}\right]
\]
si uncis in eadem significatione utimur, qua iam alibi passim usi sumus (Conf. \emph{Theorematis arithm. dem. nova} art.~4 et \emph{Theorematis fund. in doctr. de residuis quadr.} etc. \emph{Algorithm. nov.} art.~3). Contra pro valoribus ipsius $x$ his $\frac12(a-b+1),\frac12(a-b+3)\dotsfour\frac12(a-1)$ sufficiet, ut ipsis $\eta$ et $\theta$ valores positivi concilientur, sive ut $y$ maior sit quam $\frac{bx}{a}$ et minor quam $\frac{p-2ax}{2b}$ sive $\frac12b+\frac{aa-2ax}{2b}$: quare pro valore tali dato ipsius $x$ aderunt numeri $x+yi$ omnino
\[
\left[\frac12b+\frac{aa-2ax}{2b}\right]-\left[\frac{bx}{a}\right]
\]
Hinc itaque colligimus, multitudinem numerorum complexus $G$ esse
\[
g=\Sigma\left[\frac{(a+b)x}{a-b}\right]+\Sigma\left[\frac12b+\frac{aa-2ax}{2b}\right]-\Sigma\left[\frac{bx}{a}\right]
\]
ubi in termino primo summatio extendenda est per omnes valores integros ipsius $x$ ab $1$ usque ad $\frac12(a-b-1)$, in secundo ab $\frac12(a-b+1)$ usque ad $\frac12(a-1)$, in tertio ab $1$ usque ad $\frac12(a-1)$.

Si characteristica $\varphi$ in eadem significatione utimur, ut loco citato (\emph{Theorematis fund.} etc. \emph{Algor. nov.} art.~3), puta ut sit
\[
\varphi(t,u)=\left[\frac{u}{t}\right]+\left[\frac{2u}{t}\right]+\left[\frac{3u}{t}\right]\dotsfour+\left[\frac{t'u}{t}\right]
\]
denotantibus $t,u$ numeros positivos quoscunque, atque $t'$ numerum $[\frac12t]$, terminus ille primus fit $=\varphi(a-b,a+b)$, tertius $=-\varphi(a,b)$; secundus vero fit
\SourceLeaf{155}%
\[
=\frac14bb+\Sigma\left[\frac{aa-2ax}{2b}\right]
\]
Sed fit, scribendo terminos inverso ordine,
\[
\begin{aligned}
\Sigma\left[\frac{aa-2ax}{2b}\right]&=\left[\frac{a}{2b}\right]+\left[\frac{3a}{2b}\right]+\left[\frac{5a}{5b}\right]+\dotsfour+\left[\frac{(b-1)a}{2b}\right]\\
&=\varphi(2b,a)-\varphi(b,a)
\end{aligned}
\]
Formula itaque nostra sequentem induit formam:
\[
g=\varphi(a-b,a+b)+\varphi(2b,a)-\varphi(a,b)-\varphi(b,a)+\frac14bb
\]

Consideremus primo terminum $\varphi(a-b,a+b)$, qui protinus transmutatur in $\varphi(a-b,2b)+1+2+3+\text{ etc. }+\frac12(a-b-1)$ sive in
\[
\varphi(a-b,2b)+\frac18((a-b)^2-1)
\]

Dein quum per theorema generale fiat $\varphi(t,u)+\varphi(u,t)=[\frac12t]\pdot[\frac12u]$, dum $t,u$ sunt integri positivi inter se primi, habemus
\[
\varphi(a-b,2b)=\frac12b(a-b-1)-\varphi(2b,a-b)
\]
adeoque
\[
\varphi(a-b,a+b)=\frac18(aa+2ab-3bb-4b-1)-\varphi(2b,a-b)
\]

Disponamus partes ipsius $\varphi(2b,a-b)$ sequenti modo
\[
\begin{aligned}
&\left[\frac{a-b}{2b}\right]+\left[\frac{3(a-b)}{2b}\right]+\left[\frac{5(a-b)}{2b}\right]+\text{ etc. }+\left[\frac{(b-1)(a-b)}{2b}\right]\\
&+\left[\frac{a-b}{b}\right]+\left[\frac{2(a-b)}{b}\right]+\left[\frac{3(a-b)}{b}\right]+\text{ etc. }+\left[\frac{\frac12b(a-b)}{b}\right]
\end{aligned}
\]
Series secunda manifesto fit
\[
\begin{aligned}
&=\varphi(b,a-b)=\varphi(b,a)-1-2-3-\text{ etc. }-\frac12b\\
&=\varphi(b,a)-\frac18(bb+2b)
\end{aligned}
\]
seriem primam ordine terminorum inverso ita exhibemus:
\[
\begin{aligned}
&\left[\frac12(a+1-b)-\frac{a}{2b}\right]+\left[\frac12(a+3-b)-\frac{3a}{2b}\right]\\
&+\left[\frac12(a+5-b)-\frac{5a}{2b}\right]+\text{etc.}+\left[\frac12(a-1)-\frac{(b-1)a}{2b}\right]
\end{aligned}
\]
quae expressio. quum denotante $t$ numerum integrum, $u$ fractum, generaliter sit $[t-u]=t-1-[u]$, mutatur in sequentem
\[
\begin{aligned}
&\frac18b(2a-4-b)-\left[\frac{a}{2b}\right]-\left[\frac{3a}{2b}\right]-\left[\frac{5a}{2b}\right]-\text{ etc. }-\left[\frac{(b-1)a}{2b}\right]\\
&=\frac18b(2a-4-b)-\varphi(2b,a)+\varphi(b,a)
\end{aligned}
\]

